% !TEX root =../ECE343_Slides_Rogers.tex

%\begin{comment}
\section{Schedule}
\begin{frame}{Schedule}
	\def\myfontsize{footnotesize}
	\setcounter{mylnum}{24}
%\def\myfontsize{normalsize}
\begin{center}
\begin{\myfontsize}
\begin{tabular}{| c | l | c | l |}
\hline
\textbf{Lesson} 		  	& \textbf{Topics}             & \pbox{5cm}{\textbf{Reading}}	& \textbf{Notes}		\\ \hline\hline
\ifnum\lessonNum=\value{mylnum} \rowcolor{green!20!white}\fi
	\themylnum  
	& BVP (Capacitor, Coax), Electric Potential	
	& 5.16
	&  LO 5.3  
\stepcounter{mylnum}\\ \hline \hline
\ifnum\lessonNum=\value{mylnum} \rowcolor{green!20!white}\fi
	\themylnum 	
	& Electric Boundary Conditions
	& 5.18-5.22	
	& LO 5.4			
\stepcounter{mylnum}\\ \hline
\ifnum\lessonNum=\value{mylnum} \rowcolor{green!20!white}\fi
	\themylnum 	
	& BC's/Capacitance			
	& 5.23-5.24
	& 	
\stepcounter{mylnum}\\ \hline
\ifnum\lessonNum=\value{mylnum} \rowcolor{green!20!white}\fi
	\themylnum 	
	& Electrostatic Lab			
	& Lab packet
	& 	
\stepcounter{mylnum}\\ \hline
\rowcolor{gray!20!white}\multicolumn{4}{| c |}{\textbf{Chapter 6: Steady Current \& Conductivity}}	\\ \hline\hline
\ifnum\lessonNum=\value{mylnum} \rowcolor{green!20!white}\fi
	\themylnum 	
	& Current density, Conductivity, Resistance
	& Ch. 6
	&   					
\stepcounter{mylnum}\\ \hline\hline
\rowcolor{gray!20!white}\multicolumn{4}{| c |}{\textbf{Chapter 7: Magnetostatics}}						\\ \hline\hline
\ifnum\lessonNum=\value{mylnum} \rowcolor{green!20!white}\fi
	\themylnum 	
	&  Magnetic Fields, Gauss' Law for Magnetism	
	& 7.1-7.3		
	&LO 7.1-7.2 \textbf{GR 2} 				
\stepcounter{mylnum}\\ \hline
\ifnum\lessonNum=\value{mylnum} \rowcolor{green!20!white}\fi
	\themylnum 	
	& Ampere's Law 	
	& 7.4-7.5, 7.9
	& LO 7.3		
\stepcounter{mylnum}\\ \hline 
\ifnum\lessonNum=\value{mylnum} \rowcolor{green!20!white}\fi
	\themylnum 	
	&  Ampere's Law Applied
	&  7.6-7.8
	& LO 7.3
\stepcounter{mylnum}\\ \hline 
\ifnum\lessonNum=\value{mylnum} \rowcolor{green!20!white}\fi
	\themylnum	
	& Magnetic Boundary Conditions	
	& 7.10-7.14	
	& LO 7.4 		
\stepcounter{mylnum}\\ \hline 
\ifnum\lessonNum=\value{mylnum} \rowcolor{green!20!white}\fi
	\themylnum  
	& Magnetostatics Lab
	& Lab packet
	&  		
\stepcounter{mylnum}\\ \hline \hline
\rowcolor{gray!20!white}\multicolumn{4}{| c |}{\textbf{Chapter 8: Time-Varying Fields}}					\\ \hline\hline
\ifnum\lessonNum=\value{mylnum} \rowcolor{green!20!white}\fi
	\themylnum 	
	& Faraday's Law, Displacement Current, Maxwell Eqns
	& 8.3, 8.8, 8.9
	& LO 8.1
\stepcounter{mylnum}\\ \hline 
\ifnum\lessonNum=\value{mylnum} \rowcolor{green!20!white}\fi
	\themylnum	
	& Faraday's Law, Displacement Current, Maxwell Eqns 
	&  8.3, 8.8, 8.9
	& LO 8.2-8.3
\stepcounter{mylnum}\\ \hline 

\rowcolor{gray!20!white}\multicolumn{4}{| c |}{\textbf{Chapter 9: Plane Waves in Lossless Media}}					\\ \hline\hline
\ifnum\lessonNum=\value{mylnum} \rowcolor{green!20!white}\fi
	\themylnum	
	& Maxwell's Equations in lossless, source-free region		
	& 9.1-9.2
	& LO 9.1	 
\stepcounter{mylnum}\\ \hline
\ifnum\lessonNum=\value{mylnum} \rowcolor{green!20!white}\fi
	\themylnum	
	& Maxwell's Equations in lossless, source-free region		
	& 9.1-9.2
	& LO 9.1	 
\stepcounter{mylnum}\\ \hline
%\ifnum\lessonNum=\value{mylnum} \rowcolor{green!20!white}\fi
%	\themylnum	
%	& Flex day
%	& 
%	& \textbf{GR 3}			 					
%\stepcounter{mylnum}\\ \hline 
\ifnum\lessonNum=\value{mylnum} \rowcolor{green!20!white}\fi
	\themylnum	
	& Uniform Plane Waves \& Polarization
	& 9.4-9.6
	& LO 9.2-9.3			
\stepcounter{mylnum}\\ \hline 
\ifnum\lessonNum=\value{mylnum} \rowcolor{green!20!white}\fi
	\themylnum	
	& Normal Incidence
	& Vol 2, 5.6-5.7
	&  LO 9.4-9.5
\stepcounter{mylnum}\\ \hline 
\ifnum\lessonNum=\value{mylnum} \rowcolor{green!20!white}\fi
	\themylnum	
	& Oblique Incidence			
	& Vol 2, 5.6-5.7
	& LO 9.5
\stepcounter{mylnum}\\ \hline 
\ifnum\lessonNum=\value{mylnum} \rowcolor{green!20!white}\fi
	\themylnum	
	& Oblique Incidence			
	& Vol 2, 5.6-5.7
	& LO 9.5
\stepcounter{mylnum}\\ \hline 
\ifnum\lessonNum=\value{mylnum} \rowcolor{green!20!white}\fi
	\themylnum	
	& Incidence Lab			
	& TBD
	& LO 9.5
\stepcounter{mylnum}\\ \hline 
%\ifnum\lessonNum=\value{mylnum} \rowcolor{green!20!white}\fi
%	\themylnum	
%	& Wrapup
%	& 
%	& 		 					
%\stepcounter{mylnum}\\ \hline 

\end{tabular}			
\end{\myfontsize}
\end{center}

\end{frame}

%%%%%%%%%%%%%%%%%%%%%%%%%%%%%%%%%%%%%%%%%%%%%%%%%%%%%
%%%%%%%%%%%%%%%%%%%%%%%%%%%%%%%%%%%%%%%%%%%%%%%%%%%%%

\def\obja{I can articulate the significance of the displacement current and calculate it, given a time-varying electric field.}
\def\objb{I can articulate the meaning of and calculate the relaxation time for a variety of materials.}
\def\objc{ I know the forms of Maxwell's equations for time-varying fields and use them to calculate electric and magnetic fields for a variety of physical systems. }
\begin{frame}{Ch 8 Learning Objectives}
	\begin{textblock}{14.5}(0.5,3.5)
	\begin{myblock}[0.95]{}
	\begin{small}
	\begin{enumerate}[{8}.1]
		\ifnum\objNum=1
			\item \alert{ \obja }
			\item \alert{ \objb }
		\else 
			\item \obja
			\item \objb
		\fi 
		\ifnum\objNum=3
			\item \alert{ \objc }
		\else 
			\item \objc
		\fi 
	\end{enumerate}
	\end{small}
	\end{myblock}
	\end{textblock}
\end{frame}

%%%%%%%%%%%%%%%%%%%%%%%%%%%%%%%%%%%%%%%%%%%%%%%%%%%%%%
%%%%%%%%%%%%%%%%%%%%%%%%%%%%%%%%%%%%%%%%%%%%%%%%%%%%%%
\section{Current Continuity}
\begin{frame}{Displacement Current}
\begin{textblock}{8}(0.5,3)

	\begin{circuitikz}[scale=0.75]
	
		\draw (0,0) to[V=\myv{dc}, invert] ++ (0,3)
					to[R=\myr{1}, f=\myi{1}] ++(3,0)
					to[R=\myr{2}, f=\myi{2}] ++(3,0)
					to[short] ++(0,-3)
					to[short] ++(-6,0);
		\draw (3,3) to[R=\myr{3}, f=\myi{3}] ++(0,-3);
%		\node[align=left] at (6,1) {KVL$\implies$ \handoutgap{0pt}{$ \myi{1} -\myi{2}-\myi{3} = 0 $}};
	\end{circuitikz}


	\begin{circuitikz}[scale=0.75]
	
		\draw (0,0) to[sV,a^=\myv{ac}, invert] ++ (0,3)
					to[R=\myr{1}] ++(3,0)
					to[R=\myr{2}] ++(3,0)
					to[short] ++(0,-3)
					to[short] ++(-6,0);
		\draw (3,3) to[R=\myr{3}] ++(0,-3);
%		\node[align=left] at (6,1) {Now $\implies$ \handoutgap{0pt}{$-\dfrac{dQ}{dt}$}};
	\end{circuitikz}
	
\end{textblock}
\begin{textblock}{8}(8,4)
\begin{small}
\visible<0|handout:1>{
\eq{
	&\text{KVL: } \myi{1} -\myi{2}-\myi{3} = 0 \\[44pt]
	&I = -\partd{Q}{t}  \\
	&\int\vec{J}\mycdot \vec{dS} = -\partd{}{t} \iiint \rho_v \;dV \\
	&\iiint \lp \div{\vec{J}} \rp dV = \iiint \lp -\partd{}{t} \rp \rho_v \;dV \\
	&\implies \div{\vec{J}} = -\partd{\rho_v}{t} 
}
} % end visible
\end{small}
\end{textblock}

\end{frame}

%%%%%%%%%%%%%%%%%%%%%%%%%%%%%%%%%%%%%%%%%%%%%%%%%%%%%%
%%%%%%%%%%%%%%%%%%%%%%%%%%%%%%%%%%%%%%%%%%%%%%%%%%%%%%
\section{Relaxation Time}
\begin{frame}{Insight \#1: Relaxation Time}
\begin{textblock}{7}(0.5,2.5)
\visible<0|handout:1>{
\alert{
\eq{
	&\div{\vec{J}}  = -\partd{\rho_v}{t} \\
	&\text{Ohm's Law } \longrightarrow \vec{J} = \sigma \vec{E} \\
	&\implies \div{\sigma\vec{E}}  = -\partd{\rho_v}{t} \\
	&\implies \div{\vec{E}}  = -\dfrac{1}{\sigma} \partd{\rho_v}{t} \\
	&\text{...but remember Gauss? } \\
	&\implies \div{\vec{E}} = \dfrac{\rho_v}{\varepsilon} \\
	&\implies \dfrac{\rho_v}{\varepsilon}  + \dfrac{1}{\sigma} \partd{\rho_v}{t} = 0    \\
	&\implies \rho_v  = \rho_0 e^{-t/\tau} \\
	&\implies - \dfrac{1}{\tau} + \dfrac{\sigma}{\varepsilon} = 0 
	\implies \tau = \dfrac{\varepsilon}{\sigma}
} % end eq
} % end alert
} % end visible

\end{textblock}
\begin{textblock}{7}(8,3)
\visible<0|handout:1>{
\alert{
\boxeq{
	\rho_v = \rho_0 e^{-t \frac{\sigma}{\varepsilon}}
} % end eq
} % end alert
} % end visible

\scalebox{0.75}{
\begin{tikzpicture}
		\begin{axis}[ 
		domain=0:4,
		xmin=0, 
		xmax=4.5, 
		ymax=1,
		ymin=-0,
%		scaled x ticks=base 10:6,
%		xtick scale label code/.code={},
%		minor x tick num = 5,
%		minor y tick num = 5,
%		ytick={-1,0,1},
		ytick={0.37},
%		xtick={-5,-4,-3,-2,-1,0,1,2,3,4,5},;
		xtick={2},
%		xticklabels={$-\frac{\tau}{2}$,$\frac{\tau}{2}$},
		xticklabels={$\tau$},
		yticklabels={$0.37$},
		xlabel={$t$},
		ylabel={$\rho_v$},
		axis y line = middle,
		axis x line = middle,
%		x label style={xshift=-0.5cm},
		x label style={yshift=-0.5cm},
%		y label style ={xshift=1cm},
		y label style ={yshift=0.5cm},
		clip=false,
		]
		\addplot[thick, blue]  {1*exp(-x/2)};
		\addplot[dashed, blue] coordinates {(0,0.37)(2,0.37)};
		\addplot[dashed, blue] coordinates {(2,0.37)(2,0)};
		\end{axis}
\end{tikzpicture}
} % end scalebox
\end{textblock}
\end{frame}

%%%%%%%%%%%%%%%%%%%%%%%%%%%%%%%%%%%%%%%%%%%%%%%%%%%%%%
%%%%%%%%%%%%%%%%%%%%%%%%%%%%%%%%%%%%%%%%%%%%%%%%%%%%%%
\begin{frame}{Relaxation Time: Examples}
\begin{textblock}{7}(0.5,3)

\includegraphics[scale=0.5]{./ECE343/polystyrene}

\vspace{12pt}

\includegraphics[scale=0.5]{./ECE343/copper}

\end{textblock}

\begin{textblock}{8}(8,2)
\begin{myblock}[0.95]{Polystyrene}
\vspace{-12pt}
\eq{
	&\sigma = 10^{-17} \text{S/m} \text{  and  } \varepsilon = 2.56\varepsilon_0 \\ 
	& \tau = \handoutgap{0pt}{\dfrac{ 2.56\varepsilon_0\text{ F/m}}{10^{-17}\text{ S/m} } = 2.3\ee{6}\text{s} }\\
	&\handoutgap{0pt}{\qquad \approx 26\text{ days} }
}
\vspace{-12pt}
\end{myblock}

%\vspace{12pt}

\begin{myblock}[0.95]{Copper}
\vspace{-12pt}
\eq{
	&\sigma = 5.8\ee{7} \text{S/m} \text{  and  } \varepsilon = \varepsilon_0 \\ 
	&\tau = \handoutgap{0pt}{\dfrac{\varepsilon_0 \text{F/m}}{5.8\ee{7}\text{S/m} } = 1.5\ee{-19}\text{s} }
}
\vspace{-12pt}
\end{myblock}

\end{textblock}

\end{frame}

%%%%%%%%%%%%%%%%%%%%%%%%%%%%%%%%%%%%%%%%%%%%%%%%%%%%%%
%%%%%%%%%%%%%%%%%%%%%%%%%%%%%%%%%%%%%%%%%%%%%%%%%%%%%%
\section{Ampere's Law}
\begin{frame}{Insight \#2: Ampere's Law}
\begin{textblock}{7}(0.5,2.5)
\visible<0|handout:1>{
\alert{
\eq{
	&\div{\vec{J}}  = -\partd{\rho_v}{t} \\
	&\text{Ampere's Law } \longrightarrow \mycurl{\vec{H}} = \vec{J} \\
	&\div{\mycurl{\vec{H}}} = \div{\vec{J}} \\
	&0 = -\partd{\rho_v}{t} \overset{\text{?}}{=} \div{\vec{J}} \\
	&\text{Maxwell to the rescue!} \\
	& \implies \vec{J} = \underbrace{\vec{J}_c}_{\text{conduction}}+ \underbrace{\vec{J}_d }_{\text{displacement}	} \\
	&\implies \div{\vec{J}} = \div{\vec{J}_c} + \div{\vec{J}_d}  = 0  \\
	&\implies \div{\vec{J}_c}  =- \div{\vec{J}_d}  = -\partd{\rho_v}{t} \\
	&\implies  \div{\vec{J}_d}  = \partd{\rho_v}{t} 
} % end eq
} % end alert
} % end visible 

\end{textblock}
\begin{textblock}{7}(8,3.5)
\visible<0|handout:1>{
\alert{
Use Gauss' Law, ($\div{\vec{E}} = \dfrac{\rho_v}{\varepsilon}$)
\eq{
	&\div{\vec{J}_d}  = \partd{}{t} \div{\varepsilon\vec{E}}  \\
	&\qquad \vec{J} = \varepsilon\partd{\vec{E}}{t} 
}
We know $\vec{J}_c = \sigma\vec{E}$, so...
} % end alert
} % end visible
\begin{myblock}[0.95]{Ampere's Law}
\vspace{-12pt}
\eq{
	\handoutgap{0pt}{\mycurl{\vec{H}} = \vec{J}_c + \vec{J}_d = \sigma\vec{E} + \varepsilon\partd{\vec{E}}{t} }
}
\vspace{-12pt}
\end{myblock}
\end{textblock}
\end{frame}


%%%%%%%%%%%%%%%%%%%%%%%%%%%%%%%%%%%%%%%%%%%%%%%%%%%%%%
%%%%%%%%%%%%%%%%%%%%%%%%%%%%%%%%%%%%%%%%%%%%%%%%%%%%%%
\section{Lenz's Law}
\begin{frame}{Lenz's Law}
\begin{textblock}{7}(0.5,3.5)


% MAGNETIC FIELD through current loop - NS in
\def\Rx{0.7}
\def\Ry{1.14}
\def\h{0.5}
\def\H{3}
\def\W{4}
\def\NB{2}
\hspace{18pt}\scalebox{0.75}{
\begin{tikzpicture}
  
  % MAGNETIC FIELD bar magnet
  \draw (0,\Ry) arc (90:270:{\Rx} and {\Ry});
  \foreach \i [evaluate={\x=(0.3-0.1*\i)*\W; \y=(0.35*\H)*\i^2/\NB; \in=180+12*\i^2; \out=1*\i^2; \f=0.4+0.02*\i;}] in {1,...,\NB}{
    \draw[BFieldLine=\f,Bcol1] (-1*\H, 0.45*\y/\H) to[out= \out,in= \in,looseness=0.8] (\x, \y);
    \draw[BFieldLine=\f,Bcol1] (-1*\H,-0.45*\y/\H) to[out=-\out,in=-\in,looseness=0.8] (\x,-\y);
  }
  \pic[rotate=-90] (M) at (-0.95*\W,0) {magnet};
  \node[Bcol1] at (0.15*\W,0.58*\H) {$\vb{B}$};
  \draw[vector] (-0.74*\W,0) --++ (0.2*\W,0) node[right,scale=0.8] {$\vb{v}$};
  
  % INDUCED MAGNETIC FIELD
  \foreach \i [evaluate={\x=(2.0-0.5*\i)*\Rx; \y=\Ry*(\i-0.5)/(\NB-0.25); \in=-10*\i^2; \out=180-\in; \f=0.57;}] in {1,...,\NB}{
    \draw[BFieldLine=\f,Bcol2] (\x, \y) to[out= \out,in= \in,looseness=0.8] (-\x, \y);
    \draw[BFieldLine=\f,Bcol2] (\x,-\y) to[out=-\out,in=-\in,looseness=0.8] (-\x,-\y);
  }
  \node[Bcol2] at (1.6*\Rx,-1.0*\Ry) {$\vb{B}_\text{ind}$};
  
  % CIRCUIT
  \draw[white,very thick]
        (0,\Ry) arc (90:-90:{\Rx} and {\Ry});
  \draw (0,\Ry) arc (90:-90:{\Rx} and {\Ry});
  \draw[current]
    (-110:{1.1*\Rx} and {1.1*\Ry}) arc (-110:-55:{1.1*\Rx} and {1.1*\Ry})
    node[midway,below right] {$I_\text{ind}$};
  
\end{tikzpicture}
} % end scalebox

\vspace{12pt}

\scalebox{0.75}{
\begin{tikzpicture}
  
  % MAGNETIC FIELD bar magnet
  \draw (0,\Ry) arc (90:270:{\Rx} and {\Ry});
  \foreach \i [evaluate={\x=(0.3-0.1*\i)*\W; \y=(0.35*\H)*\i^2/\NB; \in=180+12*\i^2; \out=1*\i^2; \f=0.4+0.02*\i;}] in {1,...,\NB}{
    \draw[BFieldLine=\f,Bcol1] (-1*\H, 0.45*\y/\H) to[out= \out,in= \in,looseness=0.8] (\x, \y);
    \draw[BFieldLine=\f,Bcol1] (-1*\H,-0.45*\y/\H) to[out=-\out,in=-\in,looseness=0.8] (\x,-\y);
  }
  \pic[rotate=-90] (M) at (-0.95*\W,0) {magnet};
  \node[Bcol1] at (0.15*\W,0.58*\H) {$\vb{B}$};
  \draw[vector] (-1.16*\W,0) --++ (-0.2*\W,0) node[left,scale=0.8] {$\vb{v}$};
  
  % INDUCED MAGNETIC FIELD
  \foreach \i [evaluate={\x=(2.0-0.5*\i)*\Rx; \y=\Ry*(\i-0.5)/(\NB-0.25); \in=-10*\i^2; \out=180-\in; \f=0.57;}] in {1,...,\NB}{
    \draw[BFieldLine=\f,Bcol2] (-\x, \y) to[out= \in,in= \out,looseness=0.8] (\x, \y);
    \draw[BFieldLine=\f,Bcol2] (-\x,-\y) to[out=-\in,in=-\out,looseness=0.8] (\x,-\y);
  }
  \node[Bcol2] at (1.6*\Rx,-1.0*\Ry) {$\vb{B}_\text{ind}$};
  
  % CIRCUIT
  \draw[white,very thick]
        (0,\Ry) arc (90:-90:{\Rx} and {\Ry});
  \draw (0,\Ry) arc (90:-90:{\Rx} and {\Ry});
  \draw[current]
    (-60:{1.1*\Rx} and {1.1*\Ry}) arc (-60:-120:{1.1*\Rx} and {1.1*\Ry})
    node[midway,below right] {$I_\text{ind}$};
  
\end{tikzpicture}
} % end scalebox
\end{textblock}

\begin{textblock}{7.5}(8.5, 4)
\begin{myblock}[0.9]{Lenz's Law}

	A current is induced by a change in the \handoutgap{0.5in}{impressed} magnetic field. The induced  \handoutgap{0.5in}{current} results from an induced magnetic field opposed to the impressed  \handoutgap{0.5in}{magnetic} field. 

\end{myblock}
\end{textblock}
\end{frame}

%%%%%%%%%%%%%%%%%%%%%%%%%%%%%%%%%%%%%%%%%%%%%%%%%%%%%%
%%%%%%%%%%%%%%%%%%%%%%%%%%%%%%%%%%%%%%%%%%%%%%%%%%%%%%
\begin{frame}{Lenz's Law}
\begin{textblock}{7}(0.5,3)


% LOOP ENTERING B FIELD
\begin{tikzpicture}
  \def\NBx{3}  % number of B field columns
  \def\NBy{5}  % number of B field rows
  \def\Lx{1.8} % B field width
  \def\Ly{3.9} % B field height
  \def\lx{2.8} % loop width
  \def\ly{2.1} % loop height
  \def\t{0.1}  % loop thickness
  
  % LOOP
  \begin{scope}[shift={(-1.8*\lx,0.18*\Ly)}]
    \draw[metal,rounded corners=1.5*\t cm,even odd rule]
      (0,0) rectangle (\lx,\ly)
      {[rounded corners=0.5*\t cm] (\t,\t) rectangle (\lx-\t,\ly-\t)};
    \draw[vector] (\lx+1.1*\t,0.6*\ly) --++ (0.3*\lx,0) node[right] {$\vb{v}_0$};
    \draw[<->] (0,-2.1*\t) --++ (\lx,0) node[midway,fill=white,inner sep=0.8] {$\ell_x$}; %{\contour{white}{$\ell_x$}};
    \draw[<->] (-2.1*\t,0) --++ (0,\ly) node[midway,fill=white,inner sep=0.8] {$\ell_y$}; %{\contour{white}{$\ell_y$}};
  \end{scope}
  
  % MAGNETIC FIELDLINES
  \draw[Bcol!30,fill=Bcol!5] (0,0) rectangle (\Lx,\Ly);
  \foreach \i [evaluate={\y=(\i-0.5)*\Ly/\NBy}] in {1,...,\NBy}{
    \foreach \j [evaluate={\x=(\j-0.5)*\Lx/\NBx}] in {1,...,\NBx}{
      \pic[scale=1] at (\x,\y) {Bout};
    }
  }
  \node[Bcol,below right] at (\Lx,\Ly) {$\vb{B}_\text{ext}$};
  
  % AXES
  \draw[<->,thick,shift={(-0.02,-0.02)}]
    (0,1.1*\Ly) node[left] {$y$} -- (0,0) -- (1.5*\Lx,0) node[below] {$x$};
  \draw[<->] (0,-2.1*\t) --++ (\Lx,0) node[midway,fill=white,inner sep=0.8] {$L$};
  
\end{tikzpicture}

\end{textblock}
\end{frame}

%%%%%%%%%%%%%%%%%%%%%%%%%%%%%%%%%%%%%%%%%%%%%%%%%%%%%%
%%%%%%%%%%%%%%%%%%%%%%%%%%%%%%%%%%%%%%%%%%%%%%%%%%%%%%
\begin{frame}{Lenz's Law}
\begin{textblock}{7}(1.5,4)


% FLUX-CHANGING LOOP
\begin{tikzpicture}
  \def\NBx{5}
  \def\NBy{4}
  \def\L{4.0}
  \def\W{2.4}
  \def\d{4.0}
  \def\t{0.1}
  \def\s{0.64}
  
  % MAGNETIC FIELDLINES
  \foreach \i [evaluate={\y=-0.15*\W+(\i-1)*1.3*\W/(\NBy-1)}] in {1,...,\NBy}{
    \foreach \j [evaluate={\x=-0.1*\L+(\j-1)*1.2*\L/(\NBx-1)}] in {1,...,\NBx}{
      \pic[scale=1] at (\x,\y) {Bin};
    }
  }
  \node[Bcol] at (1.2*\L,1.15*\W) {$\vb{B}$};
  
  % WIRE
  \draw[metal] (\L,\t/2) -- (\t,\t/2) arc (-90:-180:\t/2) -- (\t/2,\W-\t) arc (180:90:\t/2)
                         -- (\L,\W-\t/2) -- (\L,\W+\t/2) -- (\t,\W+\t/2) arc (90:180:1.5*\t)
                         -- (-\t/2,\t) arc (-180:-90:1.5*\t) -- (\L,-\t/2) -- cycle;
  \draw[metal] (\s*\L-2.5*\t,-2*\t) rectangle++ (5*\t,\W+4*\t);
  \draw[vector] (\s*\L+1.5*\t,\W/2) --++ (0.2*\L,0) node[right] {$\vb{v}$};
  \draw[current] (\s*\L,0.3*\W) --++ (0,0.4*\W) node[above] {$I_\text{ind}$};
  \draw[current] (2*\t,0.7*\W) --++ (0,-0.4*\W) node[midway,right] {$I_\text{ind}$};
  \draw[<->] (\t/2,2*\t) -- (\s*\L-2.5*\t,2*\t) node[midway] {\contour{white}{$x$}};
  \draw[<->] (-2.1*\t,\t/2) -- (-2.1*\t,\W-\t/2) node[midway] {\contour{white}{$\ell$}};
  
\end{tikzpicture}

\end{textblock}
\end{frame}

%%%%%%%%%%%%%%%%%%%%%%%%%%%%%%%%%%%%%%%%%%%%%%%%%%%%%%
%%%%%%%%%%%%%%%%%%%%%%%%%%%%%%%%%%%%%%%%%%%%%%%%%%%%%%
\section{Faraday}
\begin{frame}{Faraday's Law}
\begin{textblock}{7}(0.5,3)
Faraday's Law is the mathematical form of Lenz's Law:
\eq{
	V = -\partd{\Psi_m}{t} = -\partd{}{t} \iint \vec{B}\mycdot\vec{dS} 
} 
Electric potential comes from the electric field: 
\visible<0|handout:1>{
\alert{
\eq{
	&V = \int \vec{E}\mycdot\vec{d\ell}  \\
	& \int \vec{E}\mycdot\vec{d\ell} = -\partd{}{t} \iint \vec{B}\mycdot\vec{dS}  \\
	&  \iint \lp \mycurl{\vec{E}} \rp \mycdot \vec{dS} =  -\partd{}{t} \iint \vec{B}\mycdot\vec{dS}  
} % end eq
} % end alert
} % end visible 
\end{textblock}

\begin{textblock}{7}(8.5,7.5)
\begin{myblock}[0.95]{Faraday's Law}

\eq{
	\handoutgap{0pt}{\mycurl{\vec{E}} = -\partd{\vec{B}}{t}} \\
	}

\end{myblock}
\end{textblock}
\end{frame}

